\documentclass[10pt,a4paper]{amsart}
\usepackage[margin=19mm]{geometry}
\usepackage{amsmath,amssymb,mathtools}
\usepackage[scr=rsfs,cal=euler]{mathalfa}
\usepackage{xfrac}
\usepackage[unicode,hidelinks,bookmarks=false]{hyperref}
\newcommand{\ie}{i.e.\ }
\newcommand{\cf}{cf.\ }
\newcommand{\resp}{respectively\ }
\newcommand{\Subsection}{\subsection}
\providecommand{\nobreakdash}{\nobreak\hbox{-}}
\setlength{\emergencystretch}{2em}
\multlinegap0pt
\usepackage{listings}
\newtheorem{theorem}{Theorem}
\newtheorem{lemma}{Lemma}
\DeclareMathOperator{\Sp}{Sp}
\newcommand{\ZZ}{\mathbb{Z}}
\newcommand{\QQ}{\mathbb{Q}}
\newcommand{\Gammaab}{\Gamma(\alpha,\beta)}
\begin{document}
\title[Certifying Arithmeticity for Two Degree-Six Symplectic Hyper]{Certifying Arithmeticity for Two Degree-Six Symplectic Hypergeometric Monodromy Groups}
\author{J. Maxwell Riestenberg, Diaaeldin Taha, Steve Trettel, Anna Wienhard}
\date{}
\maketitle
\noindent\emph{Source and licence.} Certifying Arithmeticity for Two Degree-Six Symplectic Hypergeometric Monodromy Groups. Version arXiv:2605.25935v1 (CC BY 4.0). This material is distributed under the Creative Commons Attribution 4.0 International licence, http://creativecommons.org/licenses/by/4.0/, as recorded in the arXiv OAI licence metadata for this exact version and shown on the abstract page https://arxiv.org/abs/2605.25935v1. Full rights notice: licenses/arxiv-2605.25935-CC-BY-4.0.txt.\par\smallskip
\noindent\textit{Editorial scope.} Counted unit: the original \texttt{theorem} environment \texttt{thm:main} at source lines 243--247 together with its complete proof at lines 249--374; the delivered document keeps source lines 210--375 so that every referenced label and every citation resolves inside it. Native editable source is the paper's own concatenated TeX; the AMS wrapper is editorial formatting only. No publisher class, upstream script or unrelated artwork is redistributed.
\par\smallskip\noindent\textit{Reference note.} This article ships no compiled bibliography (biblatex); the entries delivered here are composed from the author's own BibTeX (.bib) fields, with no field invented and no entry dropped.
\begin{lemma}[Bajpai--Dona--Nitsche]\label{lem:bdn}
Let $n\geq 2$, let $\Omega$ be a nondegenerate symplectic form on $\QQ^{2n}$
which is integral on $\ZZ^{2n}$, and let
$\Gamma<\Sp_{\Omega}(\ZZ)$ be Zariski-dense. Then $\Gamma$ has finite index
in $\Sp_{\Omega}(\ZZ)$ if and only if $\Gamma$ contains two transvections
\[
  X_i=1+\lambda_i x_i\Omega(x_i,\cdot),\qquad
  \lambda_i\in\QQ^\times,\quad i=1,2,
\]
whose directions $x_1,x_2$ are linearly independent and
$\Omega$-orthogonal, i.e. $\Omega(x_1,x_2)=0$.
\end{lemma}

For the groups considered here, put
\[
  T=A^{-1}B.
\]
This is the standard rank-one unipotent element associated with the local
monodromy at $1$.  For a rank-one unipotent $U$, the one-dimensional subspace
$\operatorname{im}(U-I)$ is its transvection direction.  Since conjugates of
transvections are transvections, any word $\gamma$ in $A^{\pm1},B^{\pm1}$
gives another transvection $\gamma T\gamma^{-1}$.  Therefore, once
preservation of the relevant integral symplectic form has been checked,
Lemma \ref{lem:bdn} reduces arithmeticity to finding a word $\gamma$ such
that the directions of $T$ and $\gamma T\gamma^{-1}$ are linearly independent
and $\Omega$-orthogonal.

\section{C-47 and C-55 are arithmetic}

We now prove that C-47 and C-55 are arithmetic by verifying explicit
certificates found with AlphaEvolve.  Once the witness words are fixed, the
verification is deterministic and uses only exact matrix arithmetic over $\QQ$.


\begin{theorem}\label{thm:main}
The hypergeometric monodromy groups labelled C-47 and C-55 in
\cite[Table 3]{BajpaiDonaNitsche2025Thin} have finite index in their ambient
integral symplectic groups.  In particular, both groups are arithmetic.
\end{theorem}

\begin{proof}
First consider C-47.  Set
\[
  A=C(f_{47}),\qquad B=C(g_{47}),\qquad T=A^{-1}B.
\]
Let $w_{47}$ be the following freely reduced word of length $93$:
\begin{lstlisting}
bbbbbaabbbbbbAAbbbbbbaabbbbbbAAbbbbbbAAbbbbbbABaBaaBBBBBBAAAbAbaaBaBaBaBaBAAAAAAABaBaB
BaBabaa
\end{lstlisting}
and let $\gamma=M(w_{47})$ be the corresponding group element.

A direct computation over $\ZZ$ gives
\[
  \operatorname{im}(T-I)=\QQ x_1,\qquad
  \operatorname{im}(\gamma T\gamma^{-1}-I)=\QQ x_2,
\]
where
\[
  x_1=(-5,-8,-10,-8,-5,0),
\]
and
\[
  x_2=(491566906334,537748595482,224774947812,
       73905511690,-18977654566,0).
\]
The vectors $x_1$ and $x_2$ are linearly independent.
The integral
symplectic form preserved by $A$ and $B$ is
\[
\Omega_{47}=\begin{pmatrix}
0&29&-50&51&-28&1\\
-29&0&29&-50&51&-28\\
50&-29&0&29&-50&51\\
-51&50&-29&0&29&-50\\
28&-51&50&-29&0&29\\
-1&28&-51&50&-29&0
\end{pmatrix}.
\]
Indeed,
\[
  A^t\Omega_{47}A=\Omega_{47},\qquad
  B^t\Omega_{47}B=\Omega_{47},
\]
and
\[
  \det(\Omega_{47})=1679616\neq 0.
\]
Moreover,
\[
  \Omega_{47}(x_1,x_2)=0.
\]
The same exact computation gives that $T$ and $\gamma T\gamma^{-1}$ are
rank-one unipotents and that
\[
  T \, \gamma T \gamma^{-1} = \gamma T \gamma^{-1} \, T . %[T,\gamma T\gamma^{-1}]=I.
\]
Thus $T$ and $\gamma T\gamma^{-1}$ are commuting transvections whose
directions are linearly independent and $\Omega_{47}$-orthogonal.

Now consider C-55.  Set
\[
  A=C(f_{55}),\qquad B=C(g_{55}),\qquad T=A^{-1}B.
\]
Let $w_{55}$ be the following freely reduced word of length $49$:
\begin{lstlisting}
baaaabaaaabaaaabaaaabaaaabaaaabaaaaaBABaBABAAbaaB
\end{lstlisting}
and let $\gamma=M(w_{55})$ be the corresponding group element.

A direct computation over $\ZZ$ gives
\[
  \operatorname{im}(T-I)=\QQ x_1,\qquad
  \operatorname{im}(\gamma T\gamma^{-1}-I)=\QQ x_2,
\]
where
\[
  x_1=(-4,1,2,1,-4,0),
\]
and
\[
  x_2=(40999920,-275447328,-132048384,
       236325024,314749968,0).
\]
The vectors $x_1$ and $x_2$ are linearly independent. 
The integral
symplectic form preserved by $A$ and $B$ is
\[
\Omega_{55}=\begin{pmatrix}
0&1&6&3&4&5\\
-1&0&1&6&3&4\\
-6&-1&0&1&6&3\\
-3&-6&-1&0&1&6\\
-4&-3&-6&-1&0&1\\
-5&-4&-3&-6&-1&0
\end{pmatrix}.
\]
Indeed,
\[
  A^t\Omega_{55}A=\Omega_{55},\qquad
  B^t\Omega_{55}B=\Omega_{55},
\]
and
\[
  \det(\Omega_{55})=4096\neq 0.
\]
Moreover,
\[
  \Omega_{55}(x_1,x_2)=0.
\]
The same exact computation gives that $T$ and $\gamma T\gamma^{-1}$ are
rank-one unipotents and that
\[
  T \, \gamma T \gamma^{-1} = \gamma T \gamma^{-1} \, T . %[T,\gamma T\gamma^{-1}]=I.
\]
Thus $T$ and $\gamma T\gamma^{-1}$ are commuting transvections whose
directions are linearly independent and $\Omega_{55}$-orthogonal.

In both cases, the displayed form is integral and nondegenerate, and the
identities above show that the corresponding group
$\Gamma=\langle A,B\rangle$ lies in $\Sp_\Omega(\ZZ)$.  The Zariski-density
hypothesis in Lemma \ref{lem:bdn} is supplied by the Beukers--Heckman
classification, as recalled above.  Lemma \ref{lem:bdn} therefore applies to
C-47 and C-55, and shows that each group has finite index in its ambient
integral symplectic group.
\end{proof}


\begin{thebibliography}{99}
\bibitem[]{BajpaiDonaNitsche2025Thin} Jitendra Bajpai; Daniele Dona; Martin Nitsche. Thin monodromy in {S}p(4) and {S}p(6). Journal of Algebra 663 (2025) 487-501
\end{thebibliography}
\end{document}
